\documentclass[12pt, twocolumn, landscape, a4paper]{article}

\usepackage[margin=1.5cm]{geometry}
\usepackage[utf8]{inputenc}
\usepackage[spanish,es-noshorthands]{babel}
\usepackage{amssymb, amsmath, amsbsy}
\usepackage{mathpazo}
\usepackage{graphicx}
\usepackage{tasks}
\usepackage{enumerate}
\usepackage[pdftex]{hyperref}
\hypersetup{pdftitle={Evaluación Censal Regional de Estudiantes 2017-I},pdfauthor={Luis Carrillo Gutiérrez}}

\renewcommand{\familydefault}{\sfdefault}
\renewcommand{\arraystretch}{1.8}
\pagestyle{empty}

\begin{document}
\subsection*{Matemática - Primer grado en Secundaria}
\begin{itemize}
\item{En la ciudad $X$, se tiene una granja de $3\,000$ gallinas ponedoras, estas tienen un rendimiento diario de $4$ huevos por cada $5$ gallinas ponedoras. ¿Cuántas docenas de huevo produce cada semana?
\begin{tasks}(3)
\task{$400$ docenas}
\task{$1\,350$ docenas}
\task{$1\,400$ docenas}
\task{$1\,320$ docenas}
\task{$2\,400$ docenas}
\end{tasks}
}
\item{En una balanza, se coloca de un lado, una pesa de \scalebox{1.2}{$3$}\;\;$1/5$ Kg, y en el otro \scalebox{1.2}{$1$}\;\;$2/5$ Kg ¿Cuánto falta para equlibrar la balanza?
\begin{tasks}(3)
\task{$\dfrac{9}{5}$ Kg}\task{$-4/5$ Kg}\task{$\dfrac{3}{5}$ Kg}\task{$9/5$ gr}\task{$\dfrac{7}{5}$ Kg}
\end{tasks}
}
\item{Un albañil nos explica que el \emph{estuco} es una combinación de yeso y agua, el tiempo que demora en secar sobre la superficie aplicada se estima en la siguiente tabla :

\begin{center}
	\begin{tabular}{|c|c|c|c|c|c|c|c|c|}
		\hline
		Tiempo (minutos) & 1 & 2 & 3 & 4 & 5 & 6 & 7 & 8 \\
		\hline
		El área para estucar ($m^2$) & 4 & 8 & 12 & 16 & 20 & 24 & 28 & 32 \\
		\hline
	\end{tabular}
\end{center}

¿Cuál es el tiempo promedio (en minutos) que se dispone para hacerlo?
\begin{tasks}(3)
\task{$10$ minutos}\task{$15$ minutos}\task{$16$ minutos}\task{$17$ minutos}\task{$18$ minutos}
\end{tasks}
}
\item{Una señora envía a su hijo al mercado, para que compre cuatro quintos de kilo de queso, si el kilo de queso cuesta $15$ {\tt {\scriptsize UM}}. ¿Cuánto cuesta los cuatro quintos de kilo de queso?
\begin{tasks}(5)
\task{3}\task{5}\task{12}\task{15}\task{24}
\end{tasks}
}
\item{Un comerciante compra al por mayor camisetas de deporte y sudaderas a $28$ {\tt {\scriptsize UM}} y $12$ {\tt {\scriptsize UM}} la unidad respectivamente. Si invirtió $868$ {\tt {\scriptsize UM}} ¿Cuántas prendas compró en total?, sabiendo que la cantidad de camisetas es la mayor posible.
\begin{tasks}(3)
\task{$7$ sudaderas}
\task{$28$ prendas}
\task{$34$ prendas}
\task{$35$ prendas}
\task{$40$ prendas}
\end{tasks}
}
\item{Si se resuelve correctamente la siguiente operación combinada : $$M = 10^{\,3} \div \Big[(10 \div 5)^{\,3} \times 4 - (13 - 8)^{\,2} + \sqrt[3]{27}\Big]^{\,2} - \sqrt{81\,}$$ ¿Qué resultado se obtiene / el valor de $M$ es?
\begin{tasks}(5)
\task{-1}\task{0}\task{1}\task{2}\task{3}
\end{tasks}
}
\end{itemize}
\end{document}